% ============================================================================
%  Chapter 8 — Conclusion and Future Work
%  Thesis: Generative AI, Privacy & Personal Data — TADJ Mohamed Walid, ESI SBA
%
%  DUAL-MODE FILE (standalone OR \input via \ThesisMaster). Sets \label{ch:conclusion}
%  (resolves the last forward-ref from Ch1's roadmap). SHORT synthesis from Ch1-7 +
%  discussion_notes.md §6 — re-derives nothing, introduces no number or claim. The
%  closer lands on the honest note the body earned: reduction (~40%) not elimination,
%  demonstration not certification, linkage the open residual. No triumphant ending.
% ============================================================================
\ifdefined\ThesisMaster\else
  \documentclass[11pt,a4paper]{report}
  \usepackage[utf8]{inputenc}
  \usepackage[T1]{fontenc}
  \usepackage{lmodern}
  \usepackage[margin=2.5cm]{geometry}
  \usepackage{booktabs}
  \usepackage{tabularx}
  \usepackage{array}
  \newcolumntype{L}{>{\raggedright\arraybackslash}X}
  \usepackage{amsmath,amssymb}
  \usepackage{enumitem}
  \usepackage{microtype}
  \usepackage[hidelinks]{hyperref}
  \begin{document}
\fi

\chapter{Conclusion and Future Work}
\label{ch:conclusion}

\section{What this thesis did}
\label{sec:concl-summary}

Enterprise retrieval-augmented generation makes a sensitive corpus conversationally accessible,
and in doing so it can disclose personal data to users and attackers who are not authorized to
see it. This thesis set out to measure that leakage rather than assert it: to establish how much
an undefended baseline leaks and which attacks cause it (Chapter~\ref{ch:baseline}), how much a
layered defense reduces it under matched conditions (Chapter~\ref{ch:defenses}), whether the
instrument used to measure that reduction can itself be trusted (Chapter~\ref{ch:defenses}),
which adversary is hardest to defend and at what cost to utility
(Chapter~\ref{ch:tradeoff}), all on a controlled synthetic corpus with formally defined
leakage metrics and a threat model fixed in advance (Chapter~\ref{ch:threat}). The five research
questions were posed as one linked argument, and they are best answered as one.

\section{The answer, in one thread}
\label{sec:concl-answer}

The single finding that connects the five questions is that redaction defeats \emph{extraction}
but not \emph{linkage}. Direct disclosure of a named entity is a discrete event a detector can
catch and a redactor can remove; re-identification by linkage, a person reassembled from
individually innocuous attributes scattered across documents, is \emph{emergent}, and that one
property explains the rest. Because it is emergent, no redaction layer removes it, no
deterministic detector sees it, and the only instrument that can see it, a language-model judge,
is precisely the one whose reliability degrades on the redacted and refused output a defense
produces. The leak hardest to prevent is therefore also the leak hardest to measure, and these
are not two problems but one (Chapter~\ref{ch:discussion}). Concretely, the full defense stack
reduces judged leakage from $28.8\%$ to $17.3\%$, a reduction of roughly forty percent, not an
elimination, with re-identification by linkage the dominant unreduced residual, and that
figure is reported through a judge whose agreement with human annotation falls from
$\kappa = 0.827$ on undefended output to $\kappa = 0.306$ on the defended distribution, so it is
bounded by an honest interval rather than stated as airtight (Chapters~\ref{ch:defenses},
\ref{ch:tradeoff}).

\section{Contribution}
\label{sec:concl-contribution}

The contribution is convergence, not invention: this thesis proposes no new attack and no new
defense, but integrates five capabilities that the surveyed literature had only pursued
separately (the formal $(L,S)$ metric pair, a profile-dependent evaluation engine, a
systematic eight-family campaign, layered defenses quantified with $\Delta S$, and a validated
language-model judge) into one framework applied to RAG (Chapters~\ref{ch:sota}--\ref{ch:tradeoff}).
Its sharpest result is methodological and is offered as an argument: that measuring a defended
RAG system with a single instrument is wrong in a direction the analyst cannot see, the
deterministic detector recording its own blindness to emergent leakage as protection. The
demonstration is on this corpus; the claim that the \emph{direction} of the error generalizes,
though not its magnitude, is argued in the discussion (Chapter~\ref{ch:discussion}). A
secondary methodological lesson is that judge reliability is distribution-dependent and that a
multi-agent arbitration panel sharing a model-family lineage is circular; a single validated,
cross-vendor-corroborated judge is the honest alternative. None of this amounts to a deployable,
certified product; it is a measurement framework and the findings it produced.

\section{Future work}
\label{sec:concl-future}

Each direction below answers a limitation named in Chapter~\ref{ch:discussion}.

\begin{itemize}[leftmargin=1.4em, itemsep=2pt]
  \item \textbf{A second annotator and a multi-vendor judge panel}, to close the
    single-annotator and single-provider exposure of the measurement instrument.
  \item \textbf{Defended-distribution recalibration as standard practice}, to lift the judge
    above the measured $\kappa = 0.306$ on the output defenses actually produce.
  \item \textbf{Validation beyond the single Enron probe}, where the open question is the
    judge's over-flagging when transferred to real prose.
  \item \textbf{Linkage-specific defenses}: entity-resolution or directory-aware controls for
    the genuine re-identification residual, and per-attribute recognizers for the
    detector-recall portion redaction cannot reach.
  \item \textbf{Per-configuration utility measurement}, to isolate the knee's marginal cost
    rather than bounding it between the two extremes.
  \item \textbf{Compositional, multi-actor threats}, since every attacker profile here is
    analyzed in isolation, the setting in which the linkage residual is likely to matter most.
\end{itemize}

\section{Closing}
\label{sec:concl-closing}

The contribution that outlasts the specific percentages is methodological. The relative findings
of this work (which layer carries the reduction, which adversary is hardest to defend, how the
residual is composed) are what transfer to a real deployment; the absolute numbers, measured on
a synthetic corpus, are not predictions for any particular organization. The durable lesson sits
underneath them: a defended retrieval system must be measured by an instrument that is itself
validated and recalibrated on the output the defense produces, because the alternative, trusting
a single detector, reports a protection it cannot see past. This thesis reduces measured leakage
by roughly forty percent rather than to zero, leaves re-identification by linkage as an open
problem, and stands as a demonstration of how to measure privacy leakage in RAG honestly, not
a certification that the problem is solved.

\ifdefined\ThesisMaster\else
  \end{document}
\fi
